\section{Future Works \& Conclusion}

\begin{frame}{Background \& Related Work}
    \begin{itemize}
        \item Foundation models continue to expand via hybrid generative-discriminative objectives (e.g., MAE 2.0, CLIP 2).
        \item LLM efficiency research (DiSC, prompt pruning) informs practical deployment.
        \item Multi-module orchestration (vision-language-action) inspires wireless FM fusion plans.
    \end{itemize}
\end{frame}

\begin{frame}{Future Research Question}
    \centering
    \Large
    \textit{How do we scale the three validated modules into an integrated, resource-aware Wireless FM ready for field deployment?}
\end{frame}

\begin{frame}{Summary of Contributions}
    \begin{enumerate}
        \item \textbf{DenoMAE:} Multimodal self-supervised encoder for RF signals (representation pillar).
        \item \textbf{Plug-and-Play:} LLM-based reasoning for open-set, interpretable AMC (reasoning pillar).
        \item \textbf{Meta-Task:} Learnable regularization for rapid adaptation (adaptation pillar).
    \end{enumerate}
\end{frame}

\begin{frame}{Future Work: DenoMAE 2.0}
    \textbf{Hybrid Generative-Discriminative Learning}
    \begin{itemize}
        \item Augment reconstruction with position-aware patch classification for spatial awareness.
        \item Incorporate spectrogram modality from CDRF to capture temporal-frequency cues.
        \item Goal: learn global + local structure, enabling downstream detection tasks.
    \end{itemize}
\end{frame}

\begin{frame}{Future Work: DiSC-AMC}
    \textbf{Token- and Parameter-Efficient LLM Reasoning}
    \begin{itemize}
        \item Discretize statistics to compact vocabularies.
        \item Dynamic exemplar pruning to keep prompts short.
        \item Constrained prompting to enforce structured outputs on edge devices.
    \end{itemize}
\end{frame}

\begin{frame}{Future Work: Synergistic Integration}
    \begin{itemize}
        \item Pipeline: DenoMAE 2.0 features $\rightarrow$ Meta-Task adapters $\rightarrow$ DiSC-AMC reasoning.
        \item Shared metadata (SNR, channel estimates) passed across modules for context.
        \item Target: deployable Wireless Physical-Layer FM with interpretable outputs.
    \end{itemize}
\end{frame}

\begin{frame}{Roadmap}
    \begin{itemize}
        \item Short term: scale datasets (CDRF expansions), run ablations on DenoMAE 2.0 objectives.
        \item Mid term: prototype DiSC-AMC on smaller LLMs, integrate with DenoMAE embeddings.
        \item Long term: system-level demos (spectrum monitoring, adaptive radar) with full modular stack.
    \end{itemize}
\end{frame}

\begin{frame}{Section Summary \& RQ Answer}
    \begin{itemize}
        \item Future work focuses on enhancing each pillar (representation, reasoning, adaptation) and fusing them.
        \item \textbf{Answer:} Scaling the validated modules through DenoMAE 2.0, DiSC-AMC, and their integration provides a practical path to a deployable wireless foundation model.
    \end{itemize}
\end{frame}

\begin{frame}{Conclusion}
    \centering
    \Large
    This dissertation provides a \textbf{blueprint} for the future of Wireless AI.
    \vspace{0.5cm}
    \normalsize
    By strategically integrating world knowledge from Vision and Language, we can break the ``Wireless AI Wall'' and achieve scalable, robust, and intelligent communications.
\end{frame}

\begin{frame}
    \centering
    \Huge
    Thank You
    \vspace{1cm}
    \normalsize
    Questions?
\end{frame}

